\documentclass[11pt]{article}
\usepackage[a4paper,margin=27mm]{geometry}
\usepackage[T1]{fontenc}
\usepackage{amsmath,amsthm}
\usepackage{newpxtext,newpxmath}
\usepackage{booktabs,microtype}
\usepackage[hidelinks]{hyperref}
\setlength{\parskip}{0.2em}
\setlength{\emergencystretch}{1.5em}
\newtheorem{proposition}{Proposition}
\newtheorem{theorem}[proposition]{Theorem}
\newcommand{\Sone}{\mathbb S^1}
\newcommand{\T}{\mathcal T}
\newcommand{\code}[1]{\texttt{\detokenize{#1}}}
\title{Computing Among Four Roots\\[0.3em]
  \large A Monodromy Model of Reversible Computation}
\author{Xiaoyang Liu\\[0.2em]
  \small University of Waterloo\\
  \small\href{mailto:x98liu@uwaterloo.ca}{x98liu@uwaterloo.ca}}
\date{\small November 2026}

\begin{document}
\maketitle
\begin{abstract}
How can carrying a root continuously give a complete reversible computation
model? We answer this question for the equation $(z^2-u)^2=\theta$ over
$|u|=2$ and $|\theta|=1$. Its four roots at a common control point form the
memory. Continuous feedback, read from the carried root, generates returning
control and writing motions. Two conditional motions and an unconditional
flip realize every permutation of the four roots. Each finite program
produces an actual base path whose lifted endpoint agrees with exact symbolic
execution. Integer winding numbers give a computable half-turn executor and
an inverse that retraces the motion. The construction and its refinement
from continuous paths to finite execution are formalized in Lean. The result
is a four-state reversible model with one geometric account of state,
control, writing, composition, and inversion.
\end{abstract}

\section{A root as a persistent state}

A root of an equation is determined locally by its control variables and its
current branch. Along a continuous control path, that branch can be carried
forward. When the controls return, the root can occupy a different branch.
This suggests a concrete computational question: can roots provide the
states, while their continuous continuation provides a complete instruction
and program semantics?

We give such a model with four states. Classical covering-space theory
supplies the lifting principle~\cite{hatcher}. Reversible conditional flips
and their composition belong to the established theory of reversible
computing~\cite{toffoli}. The contribution here is their explicit realization
in a continuous root space, together with a Lean refinement connecting
finite programs, generated paths, and exact execution.

Let $\Sone=\{q\in\mathbb C:|q|=1\}$. Choose the base
\[
 B=\{(u,\theta)\in\mathbb C^2:|u|=2,\ |\theta|=1\}
\]
and take the whole solution space
\begin{equation}\label{eq:space}
 E=\{(u,\theta,z):(u,\theta)\in B,\ (z^2-u)^2=\theta\}.
\end{equation}
The projection $p:E\to B$ forgets the root. The fiber over the common control
point $(2,1)$ is the memory:
\[
 M=p^{-1}(2,1)\simeq\{\sqrt3,-\sqrt3,1,-1\}.
\]
A state always includes its carried root. Returning the control point to
$(2,1)$ makes the result available as the input to the next instruction.

\begin{proposition}\label{prop:cover}
The projection $p$ is a four-sheeted covering. A global coordinate system
$H:\Sone\times\Sone\to E$ is
\begin{equation}\label{eq:coordinates}
 H(a,b)=\left(2a^2,b^2,
   \sqrt2\,a\sqrt{1+\frac{b}{2a^2}}\right),
\end{equation}
where the square root is principal.
\end{proposition}
\begin{proof}
The quantity $b/(2a^2)$ has norm $1/2$. The argument of the square root
therefore has real part at least $1/2$, so the displayed root is continuous.
Its square is $2a^2+b$, which verifies the equation defining $E$.

For a given point $(u,\theta,z)\in E$, put $b=z^2-u$. Then $b^2=\theta$
and $|b|=1$. Choose $a\in\Sone$ with $2a^2=u$. The root in
\eqref{eq:coordinates} has square $z^2$; choosing the sign of $a$ makes that
root equal to $z$. This proves surjectivity. Moreover, $b$ is recovered from
$u,z$, and
\begin{equation}\label{eq:observe}
 a=\frac{z}{\sqrt2\sqrt{1+b/u}}
\end{equation}
recovers $a$, proving injectivity. The continuous bijection $H$ is a
homeomorphism because its domain is compact and $E$ is Hausdorff.
In these coordinates $p$ is the product of two squaring maps,
$(a,b)\mapsto(2a^2,b^2)$. Each circle squaring map is a two-sheeted covering.
\end{proof}

For a path $\gamma:x\to y$ and a state in $p^{-1}(x)$, covering-space
lifting gives a unique continuous path above $\gamma$ with that initial
state~\cite[\S1.3]{hatcher}. Write $\T_\gamma$ for its endpoint map. With
$\gamma*\delta$ meaning that $\gamma$ is followed by $\delta$,
\begin{equation}\label{eq:transport}
 \T_{\gamma*\delta}=\T_\delta\circ\T_\gamma,
 \qquad \T_{\bar\gamma}=\T_\gamma^{-1}.
\end{equation}
Endpoint-fixed homotopies preserve this transport. The retained branch is
thus a memory of the path's action.

\section{One continuation rule}

The two phases in \eqref{eq:coordinates} are observations of the current
state: $b=z^2-u$, and $a$ is given by \eqref{eq:observe}. They provide the
feedback quantities used below.

Every state is regular. Indeed, $|b|=1$, and $z=0$ would imply
$|u|=|b|$, contradicting $|u|=2$. Differentiating the defining equation
therefore gives a single continuation rule,
\begin{equation}\label{eq:velocity}
 \dot z=\frac{2(z^2-u)\dot u+\dot\theta}{4z(z^2-u)}.
\end{equation}
It applies to control motion, parameter motion, and their combinations.

The formalization establishes both directions of this description. For a
real-time curve in $E$, continuity of the carried state and differentiability
of $u,\theta$ force $z$ to be differentiable and to satisfy
\eqref{eq:velocity}. The proof uses the two square relations
$b^2=\theta$ and $z^2=u+b$. Conversely, take continuous real-time functions
$u,\theta,z$ satisfying \eqref{eq:velocity} on $[0,1]$, with the stated
derivatives, a nonzero denominator, and a valid initial root. The derivative
of $(z^2-u)^2-\theta$ is zero throughout the interval, so the root constraint
is retained. When the controls lie in $B$, lift uniqueness identifies the
endpoint with $\T_\gamma$. The following feedback motions supply explicit
trajectories for this rule.

\section{State-dependent returning motions}

To let one phase select the motion of the other, choose the continuous weight
\[
 r(q)=\frac{1+\operatorname{Re}q}{2}
     =\frac{|1+q|^2}{4},\qquad q\in\Sone.
\]
It takes values in $[0,1]$, with $r(1)=1$ and $r(-1)=0$. Define
\begin{align}
 C_t(a,b)&=\bigl(ae^{\mathrm i t r(b)/2},b\bigr),\label{eq:C}\\
 W_t(a,b)&=\bigl(a,be^{\mathrm i t r(a)/2}\bigr).\label{eq:W}
\end{align}
Each is a continuous flow for all real time: its value at zero is the
identity, and motion for times $s$ and $t$ composes to motion for $s+t$.
The selecting phase remains constant during its instruction, which proves
the flow law directly. Through $H$, these are motions of actual roots in
$E$. Their control velocities are, respectively,
\[
 (\dot u,\dot\theta)=(\mathrm i u r(b),0),
 \qquad (\dot u,\dot\theta)=(0,\mathrm i\theta r(a)),
\]
and the root follows \eqref{eq:velocity}. The weight and these feedback laws
are specified data of the model.

On $M$, both phases are signs. At time $2\pi$, $C_t$ flips $a$ exactly when
$b=+1$, and $W_t$ flips $b$ exactly when $a=+1$. In every case the base
returns to $(2,1)$. We use $C=C_{2\pi}|_M$ as conditional control and
$W=W_{2\pi}|_M$ as conditional writing. Retain also the unconditional flip
$A$, obtained at time $2\pi$ from
\[
 A_t(a,b)=\bigl(ae^{\mathrm i t\operatorname{Re}b/2},b\bigr).
\]
Thus all three instructions are returning continuous feedback motions.

Name the four sign states
$p=(+,+)$, $q=(-,+)$, $r=(+,-)$, and $s=(-,-)$. Their actions are
\begin{equation}\label{eq:gates}
 C=(p\ q),\qquad W=(p\ r),\qquad A=(p\ q)(r\ s).
\end{equation}
These are standard reversible conditional flips at the level of their bit
actions~\cite{toffoli}; here each action has a specified root trajectory.

\begin{proposition}\label{prop:gates}
The conditional instructions $C,W$ generate the symmetric group on
$\{p,q,r\}$ and fix $s$. The three instructions $A,C,W$ generate $S_4$ on $M$.
\end{proposition}
\begin{proof}
The transpositions $(p\ q)$, $(p\ r)$ and their conjugate $(q\ r)$ generate
the six permutations of $\{p,q,r\}$. For the second assertion, using
right-to-left composition of permutations, put
$\sigma=CW(AC)=(p\ r\ s\ q)$. The transposition $W=(p\ \sigma(p))$
is adjacent in this four-cycle. Conjugating it by powers of $\sigma$ gives
the neighboring transpositions around the cycle; their compositions give
every transposition and hence every permutation of $M$.
\end{proof}

Uniform open-loop monodromy assigns one base path to the whole fiber.
The instruction interface here assigns a path using the current root.
Proposition~\ref{prop:gates} concerns this state-dependent feedback interface
on the same covering $p:E\to B$.

\section{Finite programs and their paths}

Encode the signs by a bijection $\beta:M\to\{0,1\}^2$, with $0$ representing
$+1$ and $1$ representing $-1$. A program is a finite word in $A,C,W$,
executed from left to right. Its symbolic operations flip the first bit
unconditionally for $A$, the first bit when the second is $0$ for $C$, and
the second bit when the first is $0$ for $W$. Applying them in sequence
defines $\operatorname{run}(P,x)$.

For an instruction $I$ and current state $m$, project its root motion over
time $[0,2\pi]$ to obtain a closed base path $\gamma_{I,m}$. For a whole
program, concatenate these paths while carrying forward each actual
intermediate state. This defines $\gamma_{P,m}$ recursively. The dependence
on $m$ is part of the program semantics.

\begin{theorem}[Program refinement]\label{thm:program}
For every finite program $P$ and every initial root $m\in M$,
\begin{equation}\label{eq:refinement}
 \beta\bigl(\T_{\gamma_{P,m}}(m)\bigr)
       =\operatorname{run}(P,\beta(m)).
\end{equation}
Reversing the instruction word restores the initial memory state. Reversing
both the word and each motion retraces the lifted computation.
\end{theorem}
\begin{proof}
Each instruction's root trajectory is a continuous lift of its projected
path, with the required initial state. Lift uniqueness identifies its
endpoint with the geometric instruction. The action in
\eqref{eq:gates} agrees with the corresponding symbolic operation.
Equation~\eqref{eq:transport} then proves \eqref{eq:refinement} by induction
on the word. Each instruction is an involution on $M$, so the reversed word
restores the initial state. The inverse flow traces each segment backward;
applying these inverse motions in reversed order returns along the original
segments.
\end{proof}

A complete program is the exchange word $CWC$. Its bit action is
$(x,y)\mapsto(y,x)$, the familiar three-conditional-flip exchange of
reversible circuit computation~\cite{toffoli}. Starting at the root $1$, whose
phases are $(+1,-1)$, its actual root states are
\begin{equation}\label{eq:example}
 1\ \xrightarrow{C}\ 1\ \xrightarrow{W}\ \sqrt3
   \ \xrightarrow{C}\ -\sqrt3.
\end{equation}
The inverse path follows these segments backward, returning to $1$.
Thus the example specifies an initial root, an executable word, every
instruction endpoint, a generated continuous path, and its inverse.

\section{Exact execution by half-turns}

For a returning instruction, the current signs $a,b\in\{+1,-1\}$ select an
integer winding pair $(k,\ell)$:
\[
\begin{array}{c|cc}
 I&k&\ell\\\hline
 A&b&0\\
 C&(1+b)/2&0\\
 W&0&(1+a)/2
\end{array}
\]
For every real normalized time $t$, define the base motion and its lifted
phases by
\begin{align}
 (u(t),\theta(t))
   &=\bigl(2e^{2\pi\mathrm i kt},e^{2\pi\mathrm i\ell t}\bigr),\label{eq:base-wind}\\
 (a(t),b(t))
   &=\bigl(ae^{\pi\mathrm i kt},be^{\pi\mathrm i\ell t}\bigr).\label{eq:lift}
\end{align}
Substitution into $H$ gives a continuous root path over
\eqref{eq:base-wind} with the selected initial root. These formulas equal
the corresponding feedback motion at physical time $2\pi t$. At $t=1$,
\begin{equation}\label{eq:half-turn}
 (a,b)\longmapsto\bigl(a(-1)^k,b(-1)^\ell\bigr).
\end{equation}
This half-turn rule is an exact executor: it follows integer windings and
retains their effect on the carried phases. Each endpoint root has the exact
name $1$, $-1$, $\sqrt3$, or $-\sqrt3$.

For inverse execution, reverse the word and negate each selected winding
pair. Formula~\eqref{eq:lift} at negative time supplies the inverse motion.
The selecting coordinate is conserved during each instruction, so its
winding selection agrees when that instruction is traversed backward.
The executor therefore reconstructs the inverse path as well as the inverse
endpoint.

The command line exposes this execution directly:
\begin{quote}\small\ttfamily
 lake exe monodromy 1 CWC\\
 lake exe monodromy --inverse -sqrt3 CWC\\
 lake exe monodromy --trace 1 CWC
\end{quote}
The first two commands realize the forward and inverse computations in
\eqref{eq:example}. The trace option reports the selected windings.
The command syntax accepts the root names \code{1}, \code{-1},
\code{sqrt3}, and \code{-sqrt3}, and words over \code{A}, \code{C}, \code{W}.
Running \code{lake exe monodromy} gives the round trip
$1\to-\sqrt3\to1$.

\section{Formalization and verification}

Lean is the native environment for the construction and its
proofs~\cite{lean4}. The development reuses mathlib's analysis, covering
spaces, permutation groups, and continuous flows~\cite{mathlib}.
The source follows the mathematical dependencies. \code{Algebra.lean} and
\code{Covering.lean} supply the common action and lifting results.
\code{Quartic.lean} constructs $E$ and its phase observations;
\code{Paths.lean} connects carried-root paths with differentiable continuation; \code{Flows.lean} defines the returning feedback
instructions; \code{Programs.lean} gives finite-word semantics and
Theorem~\ref{thm:program}. The computable definitions are collected in
\code{Execution.lean}, and \code{Lifting.lean} connects their integer
half-turn execution to the actual covering-space transport.

In particular, \code{path_correct} and \code{bits_runGeometry} connect
program paths, root states, and symbolic readout.
\code{follow_transport} proves the exact half-turn rule for arbitrary
integer winding pairs. The all-time phase identity and the inverse-segment
identity connect the executor's forward and inverse paths to the specified
continuous feedback motions. Initialization and readout use $\beta^{-1}$
and $\beta$ throughout.

The accompanying source is fixed to Lean 4.35.0 and the mathlib revision
\begin{quote}\small\ttfamily
56a005855ea17e0b0958f73ec801ae140b0a1679
\end{quote}
and can be checked from the project root with
\begin{quote}\small\ttfamily
 lake build\\
 lake env lean src/Audit.lean
\end{quote}
The axiom audit permits \code{propext}, \code{Quot.sound}, and
\code{Classical.choice}. Proofs are checked through Lean's kernel with these
standard dependencies. The verification scope comprises the formal
mathematical definitions, geometric refinement, and pure executor. Parsing,
rendering, compilation, and the execution platform supply the surrounding
command-line environment.

The model has precisely four memory states, specified feedback laws, and
finite instruction words. Its expressive result is the complete permutation
group of those states. Its continuous realization and exact execution share
the same lifting semantics: a root persists while control and writing paths
move it, and the resulting root becomes the next state.

\section*{AI contributions}
\begingroup\small
\input{../AI_CONTRIBUTIONS_STATEMENTS.txt}
\endgroup

\section*{Openness}
The manuscript and the project's original code are released under The
Unlicense. \emph{No rights reserved.} Cited works and third-party dependencies
retain their respective terms.

\begin{thebibliography}{4}\small
\bibitem{hatcher}
Allen Hatcher. \emph{Algebraic Topology}. Cambridge University Press, 2002,
\S1.3. \url{https://pi.math.cornell.edu/~hatcher/AT/ATpage.html}.

\bibitem{toffoli}
Tommaso Toffoli. \emph{Reversible Computing}. Technical Memorandum
MIT/LCS/TM-151, Massachusetts Institute of Technology, February 1980.
\url{https://hdl.handle.net/1721.1/148978}.

\bibitem{lean4}
Leonardo de Moura and Sebastian Ullrich.
\emph{The Lean 4 Theorem Prover and Programming Language}.
In \emph{Automated Deduction---CADE 28}, LNCS 12699, pp.~625--635,
Springer, 2021. \url{https://doi.org/10.1007/978-3-030-79876-5_37}.

\bibitem{mathlib}
The mathlib Community. \emph{The Lean Mathematical Library}.
In \emph{Proceedings of the 9th ACM SIGPLAN International Conference on
Certified Programs and Proofs (CPP 2020)}, pp.~367--381, ACM, 2020.
\url{https://doi.org/10.1145/3372885.3373824}.
\end{thebibliography}
\end{document}
